\subsection{BFS Multicapa sobre Grafo Implícito de Transiciones}

\graybold{Preguntas guía} - ¿cuándo usarlo?

\begin{itemize}
\item ¿El estado se mueve aplicando bloques completos de movimientos (no movimientos individuales), y cada bloque cuenta como un solo paso?
\item ¿El número de "bloques" disponibles es grande, pero muchos podrían ser idénticos entre sí?
\item ¿Todas las transiciones cuestan lo mismo (peso 1), lo que reduce el problema a BFS en vez de Dijkstra con heap?
\end{itemize}

\graybold{¿Para qué sirve?} Encontrar el camino más corto (en número de "movidas compuestas", no de instrucciones individuales) en un grafo implícito donde cada nodo es una celda del tablero y cada arista corresponde a aplicar una orden completa. Como todas las aristas cuestan lo mismo, BFS resuelve el problema sin necesidad de una cola de prioridad.

\graybold{Complejidad:} \bigO{(N \cdot M) \cdot U + \text{suma de longitudes de ordenes UNICAS}}, donde U es el número de órdenes textualmente distintas (U puede ser mucho menor que K si hay repetidas).

\graybold{Fundamento teórico:} Aunque el problema se describió como "Dijkstra", Dijkstra con TODOS los pesos iguales a 1 es exactamente BFS: la cola de prioridad se vuelve innecesaria porque nunca hay que reordenar nada. El reto real está en construir el grafo implícito eficientemente: simular cada orden completa sobre TODAS las celdas del tablero simultáneamente (vectorizado con numpy, en vez de celda por celda), y deduplicar órdenes idénticas antes de simular — si hay muchas órdenes cortas repetidas, el número de simulaciones realmente necesarias puede ser mucho menor que K.

\graybold{Ejercicio aplicado}

Un robot en una grilla N×M debe llegar de A a B. Solo puede moverse ejecutando "órdenes" completas (secuencias predefinidas de movimientos U/D/L/R), donde un movimiento inválido (muro o borde) simplemente no mueve al robot pero la orden continúa. Cada orden completa cuenta como un movimiento. Hallar el mínimo número de órdenes para llegar a B, o -1 si es imposible.

\graybold{I/O:} N,M ≤ 100 (hasta 10\textsuperscript{4} celdas), K ≤ 10\textsuperscript{4} órdenes con suma de longitudes ≤ 10\textsuperscript{4} → lectura total con tokenización (patrón 2). Es clave deduplicar órdenes repetidas por contenido antes de simular: en el peor caso (K=10\textsuperscript{4} con longitud total 10\textsuperscript{4}), la mayoría de las órdenes son forzosamente de longitud 1 (solo hay 4 posibles), colapsando a unas pocas simulaciones reales. Verificado: caso adversario de 1663 órdenes únicas corre en \textasciitilde{}1.9s.

\graybold{Código solución}

\begin{lstlisting}[language=Python]
import sys
import numpy as np
from collections import deque

def main():
    data = sys.stdin.buffer.read().split()
    idx = 0
    n = int(data[idx]); m = int(data[idx + 1]); idx += 2
    grid_rows = [data[idx + i].decode() for i in range(n)]; idx += n
    ax = int(data[idx]); ay = int(data[idx + 1]); idx += 2
    bx = int(data[idx]); by = int(data[idx + 1]); idx += 2
    k = int(data[idx]); idx += 1
    orders = [data[idx + i].decode() for i in range(k)]; idx += k

    W = m + 2   # ancho con un borde de "muro" a cada lado (evita chequear limites aparte)
    free = np.zeros((n + 2) * W, dtype=bool)
    for i in range(n):
        row = grid_rows[i]
        for j in range(m):
            free[(i + 1) * W + (j + 1)] = (row[j] == '0')

    xs, ys = np.meshgrid(np.arange(n), np.arange(m), indexing='ij')
    pos0 = ((xs + 1) * W + (ys + 1)).ravel().astype(np.int64)  # indice lineal de cada celda

    DELTA_LIN = {'U': -W, 'D': W, 'L': -1, 'R': 1}

    def simulate_order(order):
        # aplica TODA la orden a las N*M celdas de inicio simultaneamente
        pos = pos0.copy()
        for ch in order:
            d = DELTA_LIN[ch]
            pos += d
            bad = ~free[pos]
            pos[bad] -= d          # revierte solo las que chocaron con muro/borde
        return pos

    # ordenes repetidas (frecuentes cuando hay muchas cortas: solo hay 4
    # ordenes distintas de longitud 1) se simulan una unica vez
    unique_orders = {}
    for o in orders:
        if o not in unique_orders:
            unique_orders[o] = len(unique_orders)

    u = len(unique_orders)
    dest = np.empty((u, n * m), dtype=np.int16)
    for o, uid in unique_orders.items():
        p = simulate_order(o)
        cx = p // W - 1
        cy = p % W - 1
        dest[uid, :] = cx * m + cy         # de indice lineal a indice de celda [0, n*m)
    dest = np.ascontiguousarray(dest.T)     # (N*M, U): fila=celda, columna=orden unica

    start = ax * m + ay
    target = bx * m + by

    # BFS (equivale a Dijkstra con todos los pesos = 1: cada orden cuesta 1
    # "movimiento" sin importar cuantas instrucciones tenga adentro)
    dist = np.full(n * m, -1, dtype=np.int32)
    dist[start] = 0
    q = deque([start])
    while q:
        cur = q.popleft()
        if cur == target:
            break
        row = dest[cur]                     # las U ordenes unicas aplicadas desde 'cur'
        mask = dist[row] == -1              # cuales llevan a celdas nuevas (vectorizado)
        if mask.any():
            for nxt in row[mask]:
                nxt = int(nxt)
                if dist[nxt] == -1:         # evita duplicados si dos ordenes coinciden en destino
                    dist[nxt] = dist[cur] + 1
                    q.append(nxt)

    print(int(dist[target]))

main()
\end{lstlisting}